
\section{The type-I and finite-type-II languages are incomparable}\label{sec:sep}
\subsection{A fixed one-bridge separator}\label{sec:shiftseparator}
Apply the complete one-bridge assembly of Section~\ref{sec:compiler}
to the fixed root presentation
\begin{equation}\label{eq:shiftroot}
 A^*A=I,\qquad AA^*=I-P_\circ,\qquad
 \sigma_t^\varphi(A)=2^{-it}A,
\end{equation}
where $P_\circ$ is a centralizer projection. Call the resulting
finite rational Boolean game $R_*$. Here the transported label list
contains just $P_\circ$. Encode its copies by reflections, introduce
a centralizer contraction $b=e_1be_0$, and impose exactly
\[
 b^*b=e_0,\qquad bb^*=e_1-(P_\circ)_1.
\]
Use the six frame refinements and the graph-label refinement for
$P_\circ$ from Section~\ref{sec:compiler}; apply the same graphs,
triangles and rational rows. Extraction is $A=v^*b$, and
$P=(P_\circ)_0$. Thus $A^*A=e_0$, $AA^*=e_0-P$, and the modular
ratio is $1/2$. Conversely every root model of
\eqref{eq:shiftroot} extends by that assembly, with $b=e_{10}\otimes A$.
This specifies the game without a chained-bridge hypothesis.

For its type-I model take $\ell^2(\mathbb N_0)$,
$A|m\rangle=|m+1\rangle$, $P_\circ=|0\rangle\langle0|$, and
\[
 h|m\rangle=2^{-m-1}|m\rangle,\qquad
 S(h)=\sum_{m\geq0}(m+1)2^{-m-1}=2\text{ bits}.
\]
The relation $hA=Ah/2$ verifies modular covariance.
The scaffold and the finite maximally entangled compiler factors
add finite entropy, so $R_*$ is perfect spatially with finite entropy.

In any perfect finite-host strategy, extraction takes place in the
nonzero finite corner $e_0Me_0$. Its isometry $A$ must be unitary:
a faithful finite trace gives $\tau(e_0-AA^*)=0$. Consequently
$AA^*=e_0$ forces $P=0$. On the other hand KMS for the original
faithful state gives
\[
 \varphi(AA^*)=\tfrac12\varphi(A^*A)=\tfrac12\varphi(e_0),
 \qquad \varphi(P)=\tfrac12\varphi(e_0)>0,
\]
a contradiction. Therefore $R_*$ has no finite host and witnesses
$\perf(\mathrm I)\not\subseteq\perf(\II_1)$.
In particular it has no finite-dimensional perfect strategy.

\subsection{What happens to the old Weyl-clock family}
The companion reduction \cite[Sections 4--7]{Spatial} uses four Fock
isometries $A_{ab}$ on words in $\{0,1\}^2$ with
\[
 A_{ab}^*A_{cd}=\delta_{ac}\delta_{bd}I,\qquad
 \sum_{ab}A_{ab}A_{ab}^*=I-P_\circ,\qquad
 h_FA_{ab}=A_{ab}h_F/16,
\]
where $h_F(w)=(3/4)16^{-|w|}$.
Depth $m$ is a $2^m$ square with coordinates
$x=\sum a_j2^{j-1}$, $y=\sum b_j2^{j-1}$.
Its centralizer coordinate unitaries increment $x,y$ cyclically;
its boundary projections select $x=0,y=0$.
The clock labels satisfy
\[
 J_{\rm W}K_{\rm W}=Z_{\rm clk}K_{\rm W}J_{\rm W},\quad
 W|_{(x,y)}=Z_{\rm clk}^y,\quad C_{\rm ctr}|_{(x,y)}=2^{-y}.
\]
The reader equations are
\begin{equation}\label{eq:oldreader}
 (W-I)P_{\rm ret}=0,\qquad
 T_{\rm inv}(W-I)=\tfrac12C_{\rm ctr}(I-P_{\rm ret}).
\end{equation}
Spatial soundness uses finite positive density eigenspaces and a
determinant to force a finite-order phase. This determinant is unavailable
in a diffuse finite factor.

Here is a normal $\II_\infty$ model of \emph{every} input of that family.
Put $\theta=(\sqrt5-1)/2$, $\zeta=e^{2\pi i\theta}$.
In the finite rotation factor
$L^\infty(\mathbb T)\rtimes_\theta\mathbb Z$, take unitaries
$J_{\rm W},K_{\rm W}$ with commutator $\zeta I$ and put
$Z_{\rm clk}=\zeta I$.
The invariant integral and Fourier coefficient at zero give a faithful
finite trace. Irrational rotation is free and ergodic, so a central
Fourier coefficient vanishes except at zero and its zero coefficient
is constant; the algebra is a factor. It is infinite-dimensional and
hence $\II_1$. It is the hyperfinite factor $R_{\rm hyp}$ by amenability of this
crossed product and injective-factor classification \cite{Connes76}.
The construction itself only needs that it is finite.

Use the tensor density $h_F\otimes1$ on
$B(\ell^2(\mathcal W))\vN R_{\rm hyp}$.
Set $W=\zeta^y$, $C_{\rm ctr}=2^{-y}$,
$P_{\rm ret}=1_{\{y=0\}}$, and
\[
 T_{\rm inv}|_{(x,y)}=
 \begin{cases}0,&y=0,\\
 2^{-y-1}(\zeta^y-1)^{-1},&y>0.
 \end{cases}
\]
We verify its norm globally, not just by a finite numerical sweep.
For an integer $q$ nearest to $y\theta$, the nonzero integer
$q^2+qy-y^2$ factors as
$(q-y\theta)(q-y\theta')$, $\theta'=-(1+\sqrt5)/2$.
Since $|q-y\theta|\leq1/2$,
\[
 \|y\theta\|\geq\frac1{\sqrt5\,y+1/2},\qquad
 |\zeta^y-1|=2\sin(\pi\|y\theta\|)
 \geq\frac4{\sqrt5\,y+1/2}>2^{-y-1}.
\]
The sine inequality uses concavity, $\sin(\pi t)\geq2t$
for $0\leq t\leq1/2$. In particular it supplies the explicit
all-integer bound $|\sin(\pi y\theta)|\geq8/(11y)$, since
$\sqrt5<9/4$ and $y\geq1$.
For the last inequality it suffices to replace $\sqrt5$ by $9/4$:
$2^{y+3}>9y/4+1/2$ for $y=1$, and doubling proves it inductively.
Thus $T_{\rm inv}$ is a contraction and
\eqref{eq:oldreader} holds. The binary coordinate recurrences
$y\mapsto2y+b$ verify the clock and counter endpoint equations.

In the companion's tableau \cite[Sections 5.3 and 7]{Spatial} a waiting front advances one cell per time step, at $x=y$ at time $y$.
The first positive-time return pulse $P_{\rm ret}$ freezes it, and the frozen coordinate becomes the unary input of the check of the hardwired predicate; the time-zero pulse is ignored. Here $P_{\rm ret}$ selects $y=0$ only, so no positive-time pulse occurs. The front never freezes, and the companion decoder never starts a checking computation, irrespective of the hardwired input. Its finite deterministic tableau relations are
satisfied by that classical nonhalting tableau in every finite square.
The coordinate, boundary and tableau projections are diagonal on word
sectors and commute with $h_F\otimes1$.
All endpoint labels have the required centralizer type.
Lemma~\ref{lem:chain} with $k=4$ applies to exactly the conditional copies and transported labels of \cite[Section 4.4 and Lemma 4.4]{Spatial} and to the endpoint rows of \cite[Sections 5--7]{Spatial}.
Its left-endpoint labels are $P_\circ,U_x,U_y,P_x,P_y,W,C_{\rm ctr},
Z_{\rm clk},J_{\rm W},K_{\rm W}$; source-only labels are
$P_{\rm ret},T_{\rm inv}$ and the tableau projections. The activity
corners and separate refinements coincide with that compiler.
Its root mass is $1/121$.
The resulting ambient host is $\II_\infty$ for every input.
Put $r_y=2^{y+1}|\zeta^y-1|$. The global minimum of
$\log_2 r_y$ is attained at $y=1$, and
\begin{equation}\label{eq:margin}
 2.89<\log_2 r_1<2.9,\qquad
 \log_2 r_1=2.898452049757\ldots.
\end{equation}
Here is a certificate, including the infinite tail. For $y\geq3$,
the preceding bound gives
$r_y\geq2^{y+3}/(\sqrt5\,y+1/2)>8$: at $y=3$ use
$\sqrt5<9/4$, and the ratio increases because
$2(\sqrt5\,y+1/2)>\sqrt5\,(y+1)+1/2$.
For $y=2$ the exact certificate described below gives $r_2>8$.
For $y=1$ it gives $1<r_1<8$ and
$r_{1,-}^{100}>2^{289}$, $r_{1,+}^{10}<2^{29}$ for rational
bounds $r_{1,-}<r_1<r_{1,+}$.
To construct them, enclose $\sqrt5$ by integer square roots at scale
$10^{80}$, and $\pi$ by
$16\arctan(1/5)-4\arctan(1/239)$ with 60 alternating terms and
the next-term remainder. Evaluate $\sin(\pi\theta)$ at the
rational endpoints using 40 alternating terms and the next term;
sine is decreasing on this interval in $(\pi/2,\pi)$.
Round outward to denominator $10^{40}$ and compare the displayed
integer powers. For $y=2$, use
$|\sin(2\pi\theta)|=\sin(\pi(2\theta-1))$; its argument is in
$(0,\pi/2)$, where sine increases, and the same series gives
$16\sin(\pi(2\theta-1))>8$.
The bundled exact-fraction routine implements these instructions.
It certifies the $y=1,2$ enclosures and the rational base and
coefficient inequality used in the tail estimate. Together with
the induction above these prove the minimum over all positive
integers; the 3999-point floating sweep is supplementary.
The weaker positive margin is all that contractivity requires.

Lemma~\ref{lem:chain}, with $k=4$, also extracts the four
isometries of every companion game in any normal commuting host.
Two isometries with orthogonal ranges cannot lie in a nonzero finite
corner. Hence every companion game has no finite host. This additional
fact is independent of the separator $R_*$ above.

\subsection{A tracial game with no spatial perfection}
For the other direction use a linear-system game with a perfect
commuting strategy and no perfect finite-dimensional strategy,
provided by \cite{SlofstraNonclosed}.
This does not assert that linear-system games are synchronous.
The solution-group criterion \cite[Theorems 4--5]{CLS} says its
distinguished central involution $j$ is nontrivial.
The central corner $(1-\lambda(j))L(\Gamma)/2$ has normalized trace
and a perfect left/right regular strategy: equation observables
commute in each constraint, square to $I$, and their product is the
specified sign because $j=-I$ in the corner.
Use their joint spectral projections on the left and the
commutant mirrors of variable observables on the right.
Consistency and each equation are then exact.
A factor fibre is a finite factor host; tensoring with $R_{\rm hyp}$ gives
$\II_1$. This is the tracial mechanism in
\cite[Proposition 7.4]{ReadType}.

Spatial perfect play implies finite-dimensional perfect play for
linear-system games by Cleve--Mittal \cite{CM} and
\cite[Theorem 5]{CLS}; Slofstra \cite{SlofstraNonclosed} uses this
implication in the nonclosure construction. We give the argument
for completeness. Each variable appearing in a constraint has an
Alice involution with a self-adjoint Bob mirror on the state vector,
by perfect consistency. Equality of those vectors makes Alice's
trace-class marginal commute with every such involution.
The within-constraint products and parity equations hold on the
marginal support, by the zero-loss equations and support separation.
A positive density eigenvalue has a nonzero finite-dimensional
eigenspace reducing all variable involutions.
Restriction to it preserves the products and all constraint equations.
Its maximally entangled vector with transposed Bob variables is a
finite-dimensional perfect strategy, a contradiction.
For a mixed spatial state first purify; the same Alice eigenprojection
argument applies. Type-I hosting decomposes into spatial factor
fibres, so it is also excluded.
Thus $\perf(\II_1)\not\subseteq\perf(\mathrm I)$.

Every type-I factor host tensored with $R_{\rm hyp}$ has type $\II_1$ in finite
dimension and type $\II_\infty$ in infinite dimension; a further
$B(\ell^2)$ amplification handles the former case.
Every $\II_1$ host similarly amplifies to $\II_\infty$.
The two witnesses just constructed prove
\[
 \perf(\mathrm I)\ \text{and}\ \perf(\II_1)
 \text{ are incomparable},\qquad
 \perf(\mathrm I),\,\perf(\II_1)
 \subsetneq\perf(\II_\infty).
\]
Equal arithmetical completeness does not imply equal languages.
